\documentclass[11pt,a4paper]{article}

\usepackage{amsmath}
\usepackage{graphicx}
\usepackage{geometry}
\usepackage{hyperref}
\usepackage{listings}
\usepackage{xcolor}

\geometry{margin=2.5cm}

\title{Portfolio Optimization using the Markowitz Model}
\author{Adam El Gbouri}
\date{March 2026}

\lstset{
language=Python,
basicstyle=\ttfamily\small,
keywordstyle=\color{blue},
commentstyle=\color{gray},
stringstyle=\color{orange},
breaklines=true
}

\begin{document}

\maketitle

\section{Introduction}

Portfolio construction is one of the central problems in quantitative finance. 
Investors seek to allocate capital across different assets in order to maximize return while controlling risk. 

In 1952, Harry Markowitz introduced \textit{Modern Portfolio Theory (MPT)}, which formalized the relationship between risk and diversification. 
The key idea is that the risk of a portfolio depends not only on the risk of individual assets, but also on the correlations between them.

The purpose of this project is to implement a complete portfolio optimization tool in Python based on the Markowitz framework. 
The program allows users to select financial assets, retrieve historical data, compute statistical metrics, simulate thousands of possible portfolios, and determine the optimal asset allocation.

The project combines financial theory, numerical optimization and data visualization into a single interactive application.

\section{General Architecture of the Project}

The program is organized into two main components.

\begin{itemize}
\item An \textbf{asset selection interface} allowing the user to choose financial instruments.
\item A \textbf{portfolio optimization engine} implementing the Markowitz model and producing visual results.
\end{itemize}

The workflow of the program is the following:

\begin{enumerate}
\item Asset selection through a graphical interface.
\item Download of historical price data.
\item Computation of asset returns and statistical properties.
\item Monte Carlo simulation of random portfolios.
\item Numerical optimization to find optimal portfolios.
\item Generation of a graphical dashboard summarizing the results.
\end{enumerate}

This modular structure separates the user interaction from the mathematical computations.

\section{Technologies and Libraries}

The project relies on several Python libraries widely used in data science and quantitative finance.

\begin{itemize}

\item \textbf{NumPy} is used for numerical operations and matrix algebra.

\item \textbf{Pandas} provides data structures for time series manipulation and statistical calculations.

\item \textbf{SciPy} is used for constrained optimization through the Sequential Least Squares Programming (SLSQP) algorithm.

\item \textbf{Matplotlib} is used to generate graphical visualizations of the results.

\item \textbf{yfinance} allows the program to download historical financial data directly from Yahoo Finance.

\item \textbf{Tkinter} is used to build the graphical user interface.

\item \textbf{financedatabase} is an open-source library that provides a large catalog of financial assets classified by exchange, asset class and industry sector.

\end{itemize}

The financedatabase library is particularly useful because it allows users to explore different types of financial instruments such as equities, ETFs, indices, cryptocurrencies and currencies.

\section{Asset Selection Interface}

The program starts with an interactive graphical interface where users can select the assets they wish to analyze.

The interface includes:

\begin{itemize}

\item A search bar to find assets by name or ticker symbol.

\item A dynamically generated list of assets retrieved from the \texttt{financedatabase} library.

\item Filtering by major international exchanges such as NASDAQ, NYSE, Euronext, Tokyo Stock Exchange and London Stock Exchange.

\item The possibility to manually add custom ticker symbols.

\item Adjustable parameters including the analysis period and the risk-free rate.

\end{itemize}

At least two assets must be selected in order to perform portfolio optimization, since diversification effects cannot be studied with a single asset.

\section{Financial Data Retrieval}

Once the assets are selected, the program downloads historical price data using the Yahoo Finance API through the \texttt{yfinance} library.

The dataset consists of daily closing prices for each asset over the selected time period. 
These prices form a time series matrix:

\[
P =
\begin{bmatrix}
P_{1,1} & P_{1,2} & \dots & P_{1,n} \\
P_{2,1} & P_{2,2} & \dots & P_{2,n} \\
\vdots & \vdots & \ddots & \vdots \\
P_{T,1} & P_{T,2} & \dots & P_{T,n}
\end{bmatrix}
\]

where \(T\) represents the number of trading days and \(n\) represents the number of assets.

From these prices, logarithmic returns are computed in order to analyze the relative change in value from one day to the next:

\[
r_t = \ln \left(\frac{P_t}{P_{t-1}}\right)
\]

Logarithmic returns are commonly used in financial modeling because they have desirable statistical properties and are additive over time.

\section{Statistical Properties of Assets}

Using the return series, the program estimates two key quantities:

\begin{itemize}

\item the \textbf{expected return} of each asset

\item the \textbf{covariance matrix} measuring the co-movement between assets

\end{itemize}

Daily statistics are annualized assuming approximately 252 trading days per year.

The vector of expected returns is given by

\[
\mu = 252 \times \text{mean}(r)
\]

while the covariance matrix is

\[
\Sigma = 252 \times \text{cov}(r)
\]

These two quantities form the basis of the Markowitz optimization framework.

\section{Portfolio Risk and Return}

Consider a portfolio composed of \(n\) assets with weights \(w = (w_1, w_2, ..., w_n)\).

The expected return of the portfolio is a weighted combination of individual asset returns:

\[
R_p = w^T \mu
\]

The portfolio risk is measured by the standard deviation of returns, which depends on the covariance matrix:

\[
\sigma_p = \sqrt{w^T \Sigma w}
\]

These two quantities represent the fundamental trade-off between return and risk.

To evaluate risk-adjusted performance, the Sharpe ratio is used:

\[
S = \frac{R_p - R_f}{\sigma_p}
\]

where \(R_f\) is the risk-free rate.

The Sharpe ratio measures how much excess return is obtained for each unit of risk.

\section{Monte Carlo Portfolio Simulation}

Before solving the optimization problem, the program explores the space of possible portfolios using a Monte Carlo simulation.

Thousands of random portfolios are generated by assigning random weights to each asset under the constraint:

\[
\sum_{i=1}^{n} w_i = 1
\]

Each simulated portfolio is evaluated in terms of its expected return, volatility and Sharpe ratio.

This simulation produces a cloud of portfolios in the risk-return space and provides a visual understanding of diversification effects.

\section{Efficient Frontier}

Among all possible portfolios, some are clearly inefficient because another portfolio exists with both higher return and lower risk.

The set of optimal portfolios forms the \textbf{Efficient Frontier}.

Mathematically, each point on the frontier is obtained by solving the following optimization problem:

\[
\min_w w^T \Sigma w
\]

subject to:

\[
w^T \mu = R_{target}
\]

By varying the target return \(R_{target}\), the entire efficient frontier can be traced.

\section{Optimal Portfolios}

Two portfolios are particularly important in the Markowitz framework.

\subsection{Minimum Variance Portfolio}

This portfolio minimizes total portfolio volatility without considering returns.

It represents the lowest achievable risk for the given set of assets.

\subsection{Tangent Portfolio}

The tangent portfolio maximizes the Sharpe ratio and lies on the Capital Market Line.

It represents the portfolio with the best risk-adjusted performance.

In this project, both portfolios are computed numerically using the SLSQP optimization algorithm from SciPy.

\section{Dashboard Visualization}

After completing the optimization, the program generates a graphical dashboard summarizing the results.

The dashboard contains several panels.

\textbf{Efficient Frontier Plot}

A scatter plot showing simulated portfolios along with the efficient frontier, the tangent portfolio and the minimum variance portfolio.

\textbf{Normalized Price Evolution}

A time series graph displaying normalized asset prices in order to compare their historical performance.

\textbf{Portfolio Allocation}

Two pie charts illustrating the composition of the tangent portfolio and the minimum variance portfolio.

\textbf{Individual Sharpe Ratios}

A bar chart showing the Sharpe ratio of each asset individually, allowing comparison between assets and the optimal portfolio.

\section{Conclusion}

This project demonstrates how classical portfolio theory can be implemented using modern data science tools.

By combining financial theory, statistical estimation, numerical optimization and data visualization, the program provides a complete framework for portfolio analysis.

Such an implementation can serve as a basis for more advanced quantitative finance models, including dynamic portfolio allocation, factor models, machine learning forecasts or risk management techniques.

\end{document}